\documentclass[10pt,a4paper]{amsart}
\usepackage[margin=19mm]{geometry}
\usepackage{amsmath,amssymb,mathtools}
\usepackage[scr=rsfs,cal=euler]{mathalfa}
\usepackage{xfrac}
\usepackage[unicode,hidelinks,bookmarks=false]{hyperref}
\newcommand{\ie}{i.e.\ }
\newcommand{\cf}{cf.\ }
\newcommand{\resp}{respectively\ }
\newcommand{\Subsection}{\subsection}
\providecommand{\nobreakdash}{\nobreak\hbox{-}}
\setlength{\emergencystretch}{2em}
\multlinegap0pt
\theoremstyle{plain}
\newtheorem{theorem}{Theorem}
\newtheorem{lemma}[theorem]{Lemma}
\newtheorem{proposition}[theorem]{Proposition}
\newtheorem{corollary}[theorem]{Corollary}
\theoremstyle{definition}
\newtheorem{definition}[theorem]{Definition}
\theoremstyle{remark}
\newtheorem{remark}{Remark}
\usepackage{enumitem}
\newcommand{\R}{\mathbb{R}}
\newcommand{\mm}{\mathsf{M}}
\newcommand{\tx}[1]{\textnormal{\texttt{#1}}}
\def\d{\,{\rm d}}
\def\dx{\,{\rm d}x}
\def\loc{\operatorname{loc}}
\def\eqn#1$$#2$${\begin{equation}\label#1#2\end{equation}}
\title[Partial regularity for Double Phase obstacle problems]{Partial regularity for Double Phase obstacle problems: the Euler-Lagrange system for constrained local minimizers (printed as Theorem 3.1)}
\author{Filomena De Filippis, Antonella Nastasi, Cintia Pacchiano Camacho}
\date{Source: arXiv:2508.10690v2 (CC BY 4.0)}
\begin{document}
\maketitle

\noindent\emph{Source and licence.} Partial regularity for Double Phase obstacle problems. Version arXiv:2508.10690v2 (CC BY 4.0), distributed under the Creative Commons Attribution 4.0 International licence, http://creativecommons.org/licenses/by/4.0/, as recorded in the arXiv licence metadata for this exact version and shown on the saved abstract page https://arxiv.org/abs/2508.10690v2. Full licence notice: licenses/arxiv-2508.10690-CC-BY-4.0.txt.\par\smallskip

\noindent\textit{Editorial scope.} Counted unit: the Euler-Lagrange system satisfied by constrained local minimizers of the double phase functional, printed in the article as Theorem 3.1 (source0.tex lines 788--797, an unlabelled \texttt{theorem} environment whose displayed identity carries the source label \texttt{equation}), delivered together with the article's complete original proof of it (source0.tex lines 798--847, one \texttt{proof} environment of 50 lines immediately adjacent to the statement, with no gap and no deferral). No mathematics is written, repaired or re-derived: statement and proof are the article's own text line for line, in the article's own notation, and no line of either is omitted, substituted or re-flowed.

Required context retained uncounted: the article's three conditions on the target set, with their labels (lines 372--376); the article's own numbered display of the functional $\mathcal{G}$ together with the sentence that introduces the functions $g_{\alpha\beta}$ and $B^{ij}$ (lines 404--408, source label \texttt{functionalG}); the article's structural and ellipticity assumptions on those functions (lines 410--416, source label \texttt{ellipticity}); and the paragraph with the numbered display that introduces the distance function, the tubular neighbourhood of $\partial\mm$, the vector field $\nu$, the field $v$ and the notation for the boundary set (lines 784--787). These four pieces are exactly what the statement invokes and what the proof uses: the statement refers to the conditions on the target, to the functional and to the ellipticity assumptions, and the proof starts from the third line of the ellipticity display and uses the notation fixed in the paragraph above it.

References. The only citation command in any retained line is the one in that last paragraph, which cites the article's reference \texttt{F1}; this package therefore carries exactly one bibliography entry, transcribed verbatim from the article's own bibliography and printed with the article's own label 33. No other citation command occurs in any retained line, and the delivered proof contains none.

Editorial formatting only, in addition: an AMS \texttt{amsart} document -- the same class the article itself uses -- that restates the macro definitions the retained lines use (the article's own \texttt{\textbackslash R}, \texttt{\textbackslash mm}, \texttt{\textbackslash tx}, \texttt{\textbackslash d}, \texttt{\textbackslash dx}, \texttt{\textbackslash loc} and \texttt{\textbackslash eqn}), loads the article's own list-formatting package for the retained list, and restates the article's own theorem-environment declarations without the article's per-section counter. Consequently the counted theorem prints here as Theorem 1, whereas the article prints it as Theorem 3.1, and the equations of the retained spans are numbered anew in order of appearance, so the article's own display numbers are not reproduced; every cross-reference inside the retained text resolves to this wrapper's numbering, and that is the only change to the numbering. No figure, no image-inclusion line and no drawing code occur in any retained line, so no artwork is delivered or needed.

Not retained: the article's main partial regularity theorems and every other statement, its preliminaries, its later sections and its bibliography except the single cited entry; none of that material is used as a step of the delivered proof, which is complete as the author wrote it. The package ships the article's concatenated source verbatim as \texttt{sources/183262/source0.tex}, and every delivered excerpt is an exact line range of that file. No mathematical statement, hypothesis, estimate or conclusion has been rewritten, repaired or added, and printed slips, if any, are preserved literally.
\begin{enumerate}[label=\tx{\alph*.}, ref=\tx{\alph*}]
    \item\label{item:a} $\mm$ is a bounded open subset of $\R^N$ with boundary $ \partial \mm \in C^3$ ;
    \item\label{item:b}$\mm$ is a bounded subregion, with boundary $\partial \mm \in C^3$, of an $\tx{m}$-dimensional submanifold $Y$ of $\R^N$ such that $\overline{\mm} \subset \text{Int}(Y)$;
    \item\label{item:c} $\mm$ is a compact submanifold of $\R^N$ without boundary;
\end{enumerate}

\begin{align}\label{functionalG}
\mathcal{G}(u,\Omega)  &:=  \int_{\Omega} (g_{\alpha \beta}(x,u) B^{ij}(x, u) D_{\alpha}u^i D_{\beta}u^j\bigr)^{\frac{p}{2}} 
+
 a(x)\bigl(g_{\alpha \beta}(x,u) B^{ij}(x, u) D_{\alpha}u^i D_{\beta}u^j\bigr)^{\frac{q}{2}} \mathrm{d}x \notag \\ & :=  \int_{\Omega} \tx{g}(x,u,Du) \mathrm{d}x,
\end{align}

\begin{equation} \label{ellipticity}
\begin{cases}
& g_{\alpha \beta} = g_{\beta \alpha}, \ B^{ij} = B^{ji} \\
    & g, B \in C^1 \\ 
    & \ell |\eta|^2 \leq g_{\alpha\beta}(x,u)\eta_\alpha \eta_\beta, \ \ell |\xi|^2 \leq  B^{ij}(x,u) \xi^i \xi^j ,
\end{cases}   
\end{equation}

Now, following the construction outlined in \cite[Theorem 2.1]{F1}, we arrive at the following theorem that is a Euler-Lagrange system for constrained local minimizers. First we analyze case \ref{item:a}, with  $\Omega$ that is a  bounded open subset of $\R^n$, and introduce some notation. We observe that since $\partial \mm \in C^3$, the distance function $d(t):= \text{dist}(t, \partial \mm)$ is regular for $t \in \overline{\mm}$ near $\partial \mm$. By negative reflection we extend $d$ to a smooth function on a tubular neighborhood $U$ of $\partial \mm$ and define the vector fields
\eqn{vf}
$$ \nu(t):= \text{grad} d(t), \ v(x,t):= B^{-1}(x,t)(\nu(t)),$$
for $x \in \overline{\Omega}$ and $t  \in U$, $B^{-1}$ denotes the inverse of $(B^{ij})_{1 \leq i, j \leq N}$. Observe that, on  $\partial \mm$ the vector field $\nu$ is just the interior unit normal vector field. In the following we denote $\{ u \in \partial \mm \} = \{x \in \Omega: u(x) \in \partial \mm \}$.

\begin{theorem}
  Let $\Omega$ be a  bounded open subset of $\R^n$ and let  $u \in W^{1,1}_{\loc}
    (\Omega, \overline{\mm})$ be a local minimizer of the functional $\mathcal{G}$ in \eqref{functionalG}, under assumption  \eqref{ellipticity}, with $\mm$ as in \ref{item:a}. Then, for all $\varphi \in (W^{1,\tx{H}}_0\cap L^{\infty})(\Omega, \R^N)$ it holds
    \begin{align} \label{equation}
    & \int_{\Omega} \left ( pg_p(x,u,Du) + qa(x)g_q(x,u,Du) \right )\left ( G^{ij}_{\alpha \beta}(x,u)D_\alpha u^i D_\beta \varphi^j\right ) \dx \notag  \\ & \qquad + \int_{\Omega} \frac{1}{2} \left ( pg_p(x,u,Du) + qa(x)g_q(x,u,Du)\right )\left ( D_{u^l}G^{ij}_{\alpha \beta} (x,u)D_\alpha u^i D_\beta u^j \varphi^l\right )   \dx \notag \\ 
& \quad = \int_{\{ u \in \partial \mm \}} \theta \frac{\nu(u)\varphi}{\nu(u) \cdot v(x,u)} \left ( pg_p(x,u,Du) + qa(x)g_q(x,u,Du)\right )\cdot \left ( G^{ij}_{\alpha \beta}(x,u)D_\alpha u^i D_\beta (v(x,u)^j) \right . \notag \\ & \qquad \left .+ \frac{1}{2}D_{u^l} G^{ij}_{\alpha \beta} (x,u) D_\alpha u^i D_\beta u^j v(x,u)^l\right )   \dx,
    \end{align}
    where $\theta: \Omega \to [0,1]$ is a Lebesgue measurable density function, and
    $$ G^{ij}_{\alpha \beta}(x,u) := g_{\alpha \beta}(x,u) B^{ij}(x,u), \ g_\tx{s} (x,u,z) := (G^{ij}_{\alpha \beta}(x,u)z_{\alpha}^i z_{\beta}^j)^{\frac{\tx{s}}{2}-1}, \ \text{for } \tx{s}=p,q.$$
\end{theorem}

\begin{proof}
We start covering $\partial \mm$ with balls $\tx{B}_k = \tx{B}_r(y_k)$, $y_k \in \partial \mm$, $k =1,\dots,\tx{L}$ such that $\tx{B}_{3r}(y_k) \Subset U$. Then, we choose a partition of the unity $\{ \varphi_k\}$ such that
$$ \text{supp}(\varphi_k) \Subset \tx{B}_{2r}(y_k), \quad \sum_{k=1}^\tx{L} \varphi_k =1 \text{ on } \bigcup_{k=1}^\tx{L} \tx{B}_k \supset \partial \mm.$$
Now, recall that from  $(\ref{ellipticity})_{3}$, for $x \in \overline{\Omega}$ and $t  \in U$, we get
$$ \nu(t)\cdot v(x,t) >0$$ and consider the vector fields $T_{k,1}, \dots, T_{k,N-1}$ with the following properties
$$ \text{supp}(T_{k,i}) \Subset \tx{B}_{3r}(y_k), \quad T_{k,i} \cdot T_{k,j} = \delta_{i,j}, \quad T_{k,i} \cdot \nu = 0 \text{ on } \tx{B}_{2r}(y_k), \qquad \text{for } i=1,\dots, N-1.$$
Then, locally we can decompose any test vector $\psi \in C^{1}_0(\Omega, \R^N)$ into the following sum
\eqn{eta}
$$ \varphi_k (u) \psi = \eta_k v(x,u) + \sum_{i=1}^{N-1} \eta_{k,i} T_{k,i} (u)$$
where
$$ \eta_k = \varphi_k (u) \psi \cdot \nu(u)/(\nu(u)\cdot v(\cdot,u)), \quad \eta_{k,i} = \varphi_k(u) T_{k,i}(u) \psi - \eta_k v(\cdot,u) \cdot T_{k,i}(u),$$
note that $\eta_k, \eta_{k,i} \in (W^{1,\tx{H}} \cap L^{\infty})(\Omega)$ have compact support in $\Omega$. Now, consider a smooth function $h_\varepsilon: [0,\infty) \to [0,1]$  so that $h_\varepsilon (s)=1$ for $0\leq s \leq \varepsilon$, $h_\varepsilon(s)=0$ for $s \geq 2\varepsilon$ and $h_\varepsilon' \leq 0$ and any $\eta \in C^{1}_0(\Omega)$ such that $\eta \geq 0$.

Let us first study the tangential contribution in \eqref{eta}. We take a generic vector field $ T : \R^N \to \R^N$ with support contained in a ball centered in $\partial \mm$ and such that $T \cdot \nu =0$, denoting with $\Phi(s,u)$ the flow of $T$, we can choose $$v_\delta := \Phi(\delta\eta h_\varepsilon(d(u)),u)$$ as admissible test function, for  $|\delta|$ small. Then, from the minimality of $u$, we get
\begin{align} \label{eq2}
    & \int_\Omega (pg_p(x,u,Du) + qa(x)g_q(x,u,Du))G^{ij}_{\alpha  \beta}(x,u)D_\alpha  u^i D_\beta (\eta h_\varepsilon (d(u))T^j(u)) \dx \notag \\ & \qquad + \int_{\Omega} \frac{1}{2} (pg_p(x,u,Du) + qa(x)g_q(x,u,Du)) D_{u^l} G^{ij}_{\alpha \beta} (x,u) D_\alpha u^i D_\beta u^j \eta h_\varepsilon (d(u)) T^l(u) \dx =0.
\end{align}
As regards the nontangential contribution in \eqref{eta}, we observe that for a small positive $\delta$ the following test function  is admissible
$$ u_\delta := u + \delta\eta h_\varepsilon(d(u))v(\cdot,u).$$
Moreover, by Riesz representation theorem there exists a Radon measure $\lambda  \geq 0$ on $\Omega$ such that
\begin{align} \label{eq1}
    \int_\Omega  \eta \d\lambda &= \int_{\Omega} (p g_p(x,u,Du)+ qa(x) g_q(x,u,Du)) G^{ij}_{\alpha \beta}(x,u)D_{\alpha}u^i D_\beta(\eta h_\varepsilon(d(u)))v(x,u)^j) \dx  \notag \\ & \quad + \int_{\Omega}\frac{1}{2}(p g_p(x,u,Du)+ qa(x) g_q(x,u,Du))D_{u^l} G^{ij}_{\alpha \beta} (x,u)D_\alpha u^i D_\beta u^j \eta h_\varepsilon (d(u)) v(x,u)^l \dx =: \tx{I} + \tx{II}.
\end{align}
Note that, for $\tilde{\varepsilon} \neq \varepsilon$ and $|\delta| \ll 1$, the variation $u+\delta\eta(h_\varepsilon (d(u))- h_{\tilde{\varepsilon}}(d(u)))v(\cdot,u)$ is  admissible,  so $\lambda$  is independent of $\varepsilon$. Now, we  observe that since $\nu(u)^i D_\alpha u^i = D_\alpha (d(u)) =0$ on $\{ u \in \partial \mm \}$ and $h_\varepsilon '(d(u)) \leq 0$, we get
$$ \tx{I} \overset{\varepsilon \to 0}{\longrightarrow} \int_{\{ u \in \partial \mm \}} (pg_p(x,u,Du)+qa(x)g_q(x,u,Du)) G^{ij}_{\alpha \beta}(x,u) D_\alpha u^i D_\beta (v(x,u)^j) \eta \dx$$
and
$$ \tx{II} \overset{\varepsilon \to 0}{\longrightarrow}  \int_{\{ u \in \partial \mm\} } \frac{1}{2}(pg_p(x,u,Du) + qa(x)g_q(x,u,Du)) D_{u^l}G^{ij}_{\alpha \beta}(x,u) D_\alpha u^i D_\beta u^j v(x,u)^l \eta \dx.$$
Therefore, Radon Nikodym theorem gives the existence of a density function $\theta: \Omega \to [0,1]$ such that $\lambda$ can be written as
\begin{align} \label{lambda}
    \lambda & =  \theta (pg_p(x,u,Du)+qa(x)g_q(x,u,Du)) \big \{ G^{ij}_{\alpha \beta} (x,u) D_\alpha u^i D_\beta (v(x,u)^j) \notag \\ & \qquad + \frac{1}{2}D_{u^l}G^{ij}_{\alpha \beta}(x,u)D_\alpha u^i D_\beta u^j v(x,u)^l \big \} \mathscr{L}^n \lfloor \{ u \in \partial \mm\},
\end{align}
where $\mathscr{L}^n$ is the $n$-dimensional Lebesgue measure, this shows that we can take $\eta \in (W^{1,\tx{H}}_0 \cap L^\infty)(\Omega)$.


Now, we apply \eqref{eq2} with
\[
T = T_{k,i}, \quad \eta = \eta_{k,i}, \quad \text{for } k = 1, \dots, \tx{L},
\]
and then sum over \(i\). Next, we apply \eqref{eq1} with \(\eta = \eta_k\) for \(k = 1, \dots, \tx{L}\). Combining these, we obtain
\begin{align*}
  \int_{\Omega} \frac{\psi \cdot \nu(u)}{\nu(u) \cdot v(x,u)} \varphi_k(u) \d\lambda    & =  \int_\Omega (pg_p(x,u,Du) + qa(x) g_q(x,u,Du)) G^{ij}_{\alpha \beta} (x,u) \\ & \qquad \cdot  D_\alpha u^i D_\beta \left [\left ( \sum_{l=1}^{N-1} \eta_{k,l}T_{k,l}(u) + \eta_k v(x,u) \right )^j h_\varepsilon (d(u)) \right ] \dx \\ & \quad + \int_{\Omega} \frac{1}{2}(pg_p(x,u,Du) + qa(x) g_q(x,u,Du))D_{u^h}G^{ij}_{\alpha \beta}(x,u) \\ & \qquad \cdot D_\alpha u^i D_\beta u^j \left ( \sum_{l=1}^{N-1} \eta_{k,l} T_{k,l} + \eta_k v(x,u)\right )^h h_\varepsilon (d(u)) \dx.
\end{align*}
Using \eqref{eta}, the equation above turns into
\begin{align*}
   & \int_{\Omega} \frac{\psi \cdot \nu(u)}{\nu(u) \cdot v(x,u)} \varphi_k(u) \d\lambda \\ & \quad = \int_\Omega (pg_p(x,u,Du) + qa(x)g_q(x,u,Du))G^{ij}_{\alpha \beta} (x,u) D_\alpha u^i D_\beta \left ( \varphi_k(u) h_\varepsilon (d(u)) \psi^j\right ) \dx \\ & \qquad + \int_{\Omega} \frac{1}{2}(pg_p(x,u,Du) + qa(x)g_q(x,u,Du))D_{u^h}G^{ij}_{\alpha \beta}(x,u) D_\alpha u^i D_\beta u^j \varphi_k(u) h_\varepsilon (d(u))  \psi^h \dx.
\end{align*}
Now, from the last displayed equation, taking the sum with respect to $k$, observing that $\sum_{k=1}^{\tx{L}} \varphi_k =1$ on the set where $h_\varepsilon \circ d \neq 0$, using \eqref{lambda} and adding the identity
$$ \frac{d}{d\delta|0}\mathcal{G}(w_\delta, \Omega) =0, \quad w_\delta := u + \delta(1-h_\varepsilon(d(u))) \psi,$$
we finally get \eqref{equation}.
\end{proof}

\begin{thebibliography}{99}
\bibitem[33]{F1} M. Fuchs, $p$-harmonic obstacle problems, I, \emph{Annali di Matematica pura ed applicata} (IV), Vol. CLVI (1990), 127-158.
\end{thebibliography}
\end{document}
